\chapter[DPO и методы выравнивания]{Прямая оптимизация предпочтений и методы выравнивания}
\label{ch:dpo}

\begin{quote}
\itshape
Простейшее решение задачи\,---\,то, которое переставляет то, что вы уже знаете, а не добавляет новые механизмы.
\end{quote}

\bigskip

Конвейер RLHF, описанный в главе~\ref{ch:reward-models}, включает четыре отдельных нейронных сети: референсную стратегию, обучаемую стратегию, модель вознаграждения и (в PPO) критика с функцией ценности. Одновременное обучение всех четырёх требует колоссальных инженерных усилий. Модель вознаграждения необходимо обучить, оценить и заморозить до начала RL. Критик должен быть инициализирован и синхронизирован с актором. Видеопамять GPU должна вмещать как минимум две полные копии модели\,---\,а нередко и больше при использовании отсечённого суррогата PPO или вычислении штрафа KL относительно референса. На практике эта сложность является существенным барьером: RLHF с PPO в масштабе GPT-4 или Claude требует сложной распределённой инфраструктуры обучения, тщательной настройки гиперпараметров на нескольких этапах и постоянного мониторинга взлома вознаграждения.

Прямая оптимизация предпочтений (Rafailov и др., 2023)\index{Direct Preference Optimisation (DPO)} задаёт радикальный вопрос: нужна ли нам вообще модель вознаграждения? Ответ, на который в главе~\ref{ch:reward-models} намекало замечание~\ref{rem:dpo-connection},\,---\,нет. Оптимальная стратегия в рамках KL-регуляризованной цели RLHF связана с функцией вознаграждения через обратимую параметризацию. Подстановка этой параметризации в модель предпочтений Брэдли\,---\,Терри приводит к тому, что intractable partition function сокращается, оставляя функцию потерь, зависящую только от стратегии и reference policy\,---\,величин, которые можно вычислить напрямую. Модель вознаграждения устраняется полностью.

В данной главе развивается теория и практика DPO и более широкого семейства \emph{прямых алгоритмов выравнивания}\index{direct alignment algorithm}, которые он вдохновил. В разделе~\ref{sec:dpo-motivation} объясняется мотивация отказа от модели вознаграждения. В разделе~\ref{sec:dpo-derivation} приводится полный вывод цели DPO из первых принципов с каждым алгебраическим шагом. В разделе~\ref{sec:dpo-algorithm} описывается практический алгоритм обучения, требования к данным и детали реализации. В разделе~\ref{sec:dpo-variants} рассматриваются основные варианты\,---\,IPO, KTO, cDPO и ORPO,\,---\,каждый из которых устраняет отдельное ограничение базового алгоритма. В разделе~\ref{sec:dpo-limitations} анализируются структурные слабости всех offline direct alignment methods, включая проблему сдвига распределения и отсутствие online exploration. В разделе~\ref{sec:online-rl-production} объясняется, почему online RL остаётся необходимым в production-системах, и описываются инженерные паттерны, делающие его практичным. В разделе~\ref{sec:alignment-comparison} приводится исчерпывающее сравнение всех основных методов выравнивания\,---\,от PPO до GRPO\,---\,по наиболее важным на практике параметрам. В разделе~\ref{sec:ch12-summary} подводятся итоги главы.

% ===========================================================
\section{От RLHF к прямому выравниванию}
\label{sec:dpo-motivation}

\subsection{Стоимость модели вознаграждения}

Стандартный конвейер RLHF состоит из трёх этапов. Сначала обучается стратегия после supervised fine-tuning (SFT), $\pi_\mathrm{SFT}$, на данных демонстраций. Затем обучается модель вознаграждения $r_\psi$ на данных предпочтений $\mathcal{D} = \{(x^{(i)}, y_w^{(i)}, y_l^{(i)})\}$ путём минимизации функции потерь Брэдли\,---\,Терри~\eqref{eq:rm-loss}. Наконец, стратегия дообучается с помощью PPO против $r_\psi$ с учётом штрафа KL относительно $\pi_\mathrm{SFT}$. Каждый этап вносит свою инженерную сложность и источники отказов.

\index{reward model!cost of training}
Этап модели вознаграждения особенно обременителен. Он требует отдельного набора гиперпараметров (скорость обучения, размер батча, margin), отдельного протокола оценки (точность попарного сравнения на отложенной выборке) и шага развёртывания, делающего обученную модель доступной в качестве сигнала вознаграждения во время RL. Модель вознаграждения, как правило, имеет тот же размер, что и стратегия, то есть должна помещаться в GPU memory вместе со стратегией, reference policy и PPO value function\,---\,итого до четырёх копий модели. При масштабе 70 миллиардов параметров это приближается к практическому пределу даже самых крупных GPU-кластеров.

Помимо памяти, модель вознаграждения является источником нестабильности. Если она недообучена, она не обеспечит полезный градиентный сигнал, и стратегия остановится. Если перетренирована, стратегия быстро обнаружит и начнёт эксплуатировать её слабые места (взлом вознаграждения). Поскольку модель вознаграждения заморожена во время PPO, у неё нет механизма адаптации по мере того, как распределение стратегии смещается от данных, на которых была обучена модель вознаграждения,\,---\,форма ковариантного сдвига, который усугубляется по мере продвижения обучения.

\subsection{Простота, память и стабильность}

Прямой алгоритм выравнивания\index{direct alignment algorithm} заменяет двухэтапный конвейер <<модель вознаграждения, затем RL>> единой целевой функцией в стиле обучения с учителем, вычисляемой непосредственно на данных предпочтений. Преимущества существенны.

\begin{description}[leftmargin=!,labelwidth=3.5cm]
  \item[Memory footprint] DPO требует только двух копий модели: обучаемой стратегии и замороженной reference policy. Это вдвое снижает требования к памяти по сравнению с PPO, которому нужны стратегия, reference policy, модель вознаграждения и критик с функцией ценности.
  \item[Стабильность обучения] Функция потерь DPO представляет собой цель бинарной кросс-энтропии. Она гладкая, имеет ограниченные градиенты и не страдает от нестабильностей, связанных с отсечённой суррогатной целью PPO или движущимся сигналом вознаграждения онлайн-RL.
  \item[Отсутствие reward hacking] Поскольку явной модели вознаграждения для переобучения нет, DPO не может демонстрировать классическую картину reward hacking. Стратегия не может эксплуатировать разрыв между proxy reward model и истинными предпочтениями человека, так как сигнал предпочтений встроен непосредственно в функцию потерь.
  \item[Простота реализации] DPO сводится примерно к двадцати строкам на PyTorch. Не нужен буфер опыта, сеть критика, оценка преимущества или отсечение PPO. Весь конвейер напоминает дообучение с учителем, что делает его доступным для практиков без инфраструктуры RL.
\end{description}

Эти преимущества имеют свою цену, которую разделы~\ref{sec:dpo-limitations} и~\ref{sec:online-rl-production} рассматривают подробно. Критическое ограничение состоит в том, что DPO является принципиально \emph{offline}-методом: он обучается на фиксированном наборе данных предпочтений, собранных при другой (обычно SFT) стратегии. По мере того как обученная стратегия удаляется от reference policy, данные предпочтений устаревают, и неявный сигнал вознаграждения перестаёт точно отражать то, что улучшенная стратегия предпочтёт на самом деле. Эта проблема сдвига распределения неизбежна в любом offline direct alignment algorithm, и она мотивирует продолжающуюся актуальность online RL, описанного в разделе~\ref{sec:online-rl-production}.

% ===========================================================
\section{Вывод DPO}
\label{sec:dpo-derivation}

\subsection{KL-ограниченная цель}

Начнём с KL-регуляризованной цели RLHF, введённой в главе~\ref{ch:reward-models}. При заданном промпте $x$ и функции вознаграждения $r$ цель имеет вид
\begin{equation}
\label{eq:dpo-rlhf-obj}
J(\pi_\theta) = \E_{x \sim \mathcal{D},\; y \sim \pi_\theta(\cdot|x)}\!\bigl[r(x, y)\bigr]
- \beta\,\KL{\pi_\theta(\cdot|x)}{\pi_\mathrm{ref}(\cdot|x)},
\end{equation}
где $\beta > 0$ управляет силой штрафа KL, а $\pi_\mathrm{ref}$\,---\,SFT reference policy. Цель уравновешивает максимизацию вознаграждения и близость к референсу. В теореме~\ref{thm:kl-optimal-policy} было показано, что единственная оптимальная стратегия равна
\begin{equation}
\label{eq:dpo-optimal-policy}
\pi^*(y \mid x) = \frac{\pi_\mathrm{ref}(y \mid x)\,\exp\!\bigl(r(x, y)/\beta\bigr)}{Z(x)},
\end{equation}
где статистическая сумма $Z(x) = \sum_y \pi_\mathrm{ref}(y \mid x)\,\exp(r(x, y)/\beta)$.

\subsection{Параметризация оптимальной стратегии}

Взяв логарифм обеих частей~\eqref{eq:dpo-optimal-policy} и перегруппировав, получаем \emph{параметризацию вознаграждения}:
\begin{equation}
\label{eq:dpo-reparam}
r(x, y) = \beta\log\frac{\pi^*(y \mid x)}{\pi_\mathrm{ref}(y \mid x)} + \beta\log Z(x).
\end{equation}
Это тождество выражает вознаграждение через оптимальную стратегию и статистическую сумму. Два свойства этого тождества принципиально важны.

Во-первых, тождество справедливо для \emph{любой} функции вознаграждения $r$ и соответствующей ей оптимальной стратегии $\pi^*$. Это не приближение; это точное алгебраическое следствие условий оптимальности.

Во-вторых, статистическая сумма $Z(x)$ зависит только от промпта $x$, но не от конкретного ответа $y$. Это означает, что при вычислении \emph{разности} вознаграждений между двумя ответами $y_w$ и $y_l$ на один и тот же промпт $x$ статистическая сумма сокращается:
\begin{equation}
\label{eq:dpo-reward-diff}
r(x, y_w) - r(x, y_l)
= \beta\log\frac{\pi^*(y_w \mid x)}{\pi_\mathrm{ref}(y_w \mid x)}
- \beta\log\frac{\pi^*(y_l \mid x)}{\pi_\mathrm{ref}(y_l \mid x)}.
\end{equation}
Неподдающийся вычислению член $\log Z(x)$ исчез.

\subsection{Подстановка в модель Брэдли\,---\,Терри}

Вспомним модель предпочтений Брэдли\,---\,Терри из главы~\ref{ch:reward-models}: вероятность того, что человек предпочтёт $y_w$ над $y_l$ при промпте $x$, равна
\begin{equation}
\label{eq:dpo-bt}
\Prob(y_w \succ y_l \mid x) = \sigma\bigl(r(x, y_w) - r(x, y_l)\bigr).
\end{equation}
Подставляя разность вознаграждений~\eqref{eq:dpo-reward-diff} в~\eqref{eq:dpo-bt} и заменяя $\pi^*$ параметрической стратегией $\pi_\theta$ (поскольку именно $\pi^*$ мы хотим выучить), получаем
\begin{equation}
\label{eq:dpo-pref-prob}
\Prob(y_w \succ y_l \mid x) = \sigma\!\left(\beta\log\frac{\pi_\theta(y_w \mid x)}{\pi_\mathrm{ref}(y_w \mid x)}
- \beta\log\frac{\pi_\theta(y_l \mid x)}{\pi_\mathrm{ref}(y_l \mid x)}\right).
\end{equation}
Это вероятность предпочтения DPO. Она зависит только от log-ratio стратегии к reference policy, которые вычисляются двумя прямыми проходами через языковую модель (один для $y_w$ и один для $y_l$).

\subsection{Функция потерь DPO}

При наличии набора данных $\mathcal{D} = \{(x^{(i)}, y_w^{(i)}, y_l^{(i)})\}_{i=1}^N$ предпочтений людей отрицательное логарифмическое правдоподобие наблюдаемых предпочтений согласно модели~\eqref{eq:dpo-pref-prob} является функцией потерь DPO:

\begin{definition}[Функция потерь DPO]
\label{def:dpo-loss}
\index{DPO loss}
\textbf{Функция потерь прямой оптимизации предпочтений (DPO)} имеет вид
\begin{equation}
\label{eq:dpo-loss}
\boxed{
\mathcal{L}_\mathrm{DPO}(\theta)
= -\frac{1}{N}\sum_{i=1}^{N}
\log\sigma\!\left(
\beta\log\frac{\pi_\theta(y_w^{(i)} \mid x^{(i)})}{\pi_\mathrm{ref}(y_w^{(i)} \mid x^{(i)})}
- \beta\log\frac{\pi_\theta(y_l^{(i)} \mid x^{(i)})}{\pi_\mathrm{ref}(y_l^{(i)} \mid x^{(i)})}
\right).
}
\end{equation}
\end{definition}

Функция потерь DPO представляет собой потерю бинарной кросс-энтропии, где «логит» для каждой пары\,---\,разность неявных логарифмических вознаграждений между предпочтительным и нежелательным ответами. Минимизация $\mathcal{L}_\mathrm{DPO}$ эквивалентна максимизации правдоподобия наблюдаемых аннотаций предпочтений в рамках параметризованной модели Брэдли\,---\,Терри.

\begin{proposition}[Структура градиента функции потерь DPO]
\label{prop:dpo-gradient}
\index{DPO loss!gradient}
Для одной пары предпочтений $(x, y_w, y_l)$ пусть
\[
h_\theta = \beta\log\frac{\pi_\theta(y_w|x)}{\pi_\mathrm{ref}(y_w|x)}
         - \beta\log\frac{\pi_\theta(y_l|x)}{\pi_\mathrm{ref}(y_l|x)}
\]
будет разностью неявных вознаграждений. Тогда градиент $\mathcal{L}_\mathrm{DPO}$ по $\theta$ равен
\begin{equation}
\label{eq:dpo-gradient}
\nabla_\theta \mathcal{L}_\mathrm{DPO}
= -\beta\,\sigma(-h_\theta)\bigl[
  \nabla_\theta\log\pi_\theta(y_w|x)
  - \nabla_\theta\log\pi_\theta(y_l|x)
\bigr].
\end{equation}
Множитель $\sigma(-h_\theta) \in (0,1)$ взвешивает обновление: он велик (близок к~$1$), когда стратегия присваивает схожее неявное вознаграждение обоим ответам (пара сложная), и мал (близок к~$0$), когда стратегия уже явно предпочитает $y_w$ (пара лёгкая). Градиент увеличивает логарифмическую вероятность $y_w$ и уменьшает логарифмическую вероятность $y_l$ с весом, пропорциональным неуверенности текущей модели.
\end{proposition}

\begin{proof}
Пусть $\mathcal{L} = -\log\sigma(h_\theta)$. Поскольку $\frac{d}{dh}\log\sigma(h) = \sigma(-h) = 1 - \sigma(h)$, имеем
\[
\nabla_\theta \mathcal{L} = -\sigma(-h_\theta)\,\nabla_\theta h_\theta.
\]
Вычислим $\nabla_\theta h_\theta$: так как референсная модель $\pi_\mathrm{ref}$ не зависит от $\theta$,
\[
\nabla_\theta h_\theta = \beta\,\nabla_\theta\log\pi_\theta(y_w|x) - \beta\,\nabla_\theta\log\pi_\theta(y_l|x).
\]
Объединяя, получаем~\eqref{eq:dpo-gradient}.
\end{proof}

\begin{remark}[DPO как дообучение с учителем с контрастным взвешиванием]
\label{rem:dpo-sft}
Уравнение~\eqref{eq:dpo-gradient} показывает, что DPO структурно идентично двум одновременным обновлениям дообучения с учителем: положительному (повышающему вероятность) обновлению для $y_w$ и отрицательному (понижающему вероятность) обновлению для $y_l$, оба масштабируются весом доверия $\sigma(-h_\theta)$ для каждой пары. Это делает DPO совместимым с теми же трюками реализации, что и SFT: контрольные точки градиента, смешанная точность, адаптеры LoRA и оптимизация DeepSpeed ZeRO применяются без изменений.
\end{remark}

\begin{example}[Вычисление функции потерь DPO]
\label{ex:dpo-computation}
Предположим, что для одной пары $(x, y_w, y_l)$ логарифмические вероятности при текущей стратегии и референсе равны:
\[
\log\pi_\theta(y_w|x) = -3{,}2, \quad \log\pi_\mathrm{ref}(y_w|x) = -3{,}5,
\quad \log\pi_\theta(y_l|x) = -2{,}8, \quad \log\pi_\mathrm{ref}(y_l|x) = -2{,}6.
\]
При $\beta = 0{,}1$ разность неявных вознаграждений равна
\[
h_\theta = 0{,}1\cdot(-3{,}2 - (-3{,}5)) - 0{,}1\cdot(-2{,}8 - (-2{,}6))
= 0{,}1\cdot(0{,}3) - 0{,}1\cdot(-0{,}2) = 0{,}03 + 0{,}02 = 0{,}05.
\]
Потеря для этой пары равна $-\log\sigma(0{,}05) = -\log(0{,}5125) \approx 0{,}669$. Вес доверия $\sigma(-0{,}05) \approx 0{,}4875$, что указывает на почти максимальную неопределённость для данной пары, и градиентное обновление будет большим. После обучения, когда $h_\theta$ достигнет~$3$, потеря станет $-\log\sigma(3) \approx 0{,}049$, а вес доверия упадёт до $\sigma(-3) \approx 0{,}047$, отражая, что стратегия научилась явно предпочитать $y_w$.
\end{example}

% ===========================================================
\section{Алгоритм DPO}
\label{sec:dpo-algorithm}

\subsection{Процедура обучения}

Процедура обучения DPO удивительно проста. Она требует двух экземпляров модели\,---\,обучаемой стратегии $\pi_\theta$ и замороженной референсной стратегии $\pi_\mathrm{ref}$\,---\,и набора данных предпочтений. Референсная стратегия\,---\,это, как правило, контрольная точка SFT, замороженная с начала обучения DPO.

\begin{algorithm}[ht]
\caption{Прямая оптимизация предпочтений (DPO)}
\label{alg:dpo}
\KwIn{Контрольная точка SFT $\pi_\mathrm{ref}$; набор данных предпочтений $\mathcal{D}$; коэффициент KL $\beta$; скорость обучения $\alpha$; размер батча $B$}
\KwOut{Выровненная стратегия $\pi_\theta$}
Инициализировать $\pi_\theta \leftarrow \pi_\mathrm{ref}$\;
Заморозить $\pi_\mathrm{ref}$ (без обновления градиентов)\;
\For{каждый мини-батч $\{(x^{(i)}, y_w^{(i)}, y_l^{(i)})\}_{i=1}^{B} \subset \mathcal{D}$}{
  \For{каждая пара $i = 1,\ldots, B$ параллельно}{
    Вычислить $\ell_w^{(i)} \leftarrow \log\pi_\theta(y_w^{(i)}|x^{(i)})$ и $\ell_l^{(i)} \leftarrow \log\pi_\theta(y_l^{(i)}|x^{(i)})$\;
    Вычислить $\hat\ell_w^{(i)} \leftarrow \log\pi_\mathrm{ref}(y_w^{(i)}|x^{(i)})$ и $\hat\ell_l^{(i)} \leftarrow \log\pi_\mathrm{ref}(y_l^{(i)}|x^{(i)})$ (без градиентов)\;
    Вычислить $h^{(i)} \leftarrow \beta(\ell_w^{(i)} - \hat\ell_w^{(i)}) - \beta(\ell_l^{(i)} - \hat\ell_l^{(i)})$\;
  }
  Вычислить $\mathcal{L} \leftarrow -\frac{1}{B}\sum_{i=1}^B \log\sigma(h^{(i)})$\;
  $\theta \leftarrow \theta - \alpha\,\nabla_\theta\,\mathcal{L}$\;
  Журнал: средний отступ вознаграждения $\frac{1}{B}\sum_i h^{(i)}$; доля правильных $\frac{1}{B}\sum_i \ind[h^{(i)} > 0]$\;
}
\Return $\pi_\theta$
\end{algorithm}

\subsection{Требования к данным}

\index{preference dataset!DPO requirements}
DPO требует набора данных предпочтений вида $\mathcal{D} = \{(x, y_w, y_l)\}$. Несколько практических соображений определяют качество и полезность этих данных.

\textbf{Разрыв в качестве ответов.} Функция потерь DPO наиболее информативна, когда между $y_w$ и $y_l$ существует явный разрыв в качестве. Если оба ответа почти идентичны, $h_\theta$ будет близок к нулю для любой разумной стратегии, и градиентное обновление будет направлено в произвольную сторону. Следует отбирать пары так, чтобы $y_w$ и $y_l$ демонстрировали существенно разный уровень качества.

\textbf{Источник ответов.} В оригинальной статье DPO оба ответа $y_w$ и $y_l$ сэмплируются из $\pi_\mathrm{ref}$ (модели SFT). Это наиболее последовательный выбор для автономного DPO: обучающие ответы взяты из того же распределения, что и референсная стратегия, поэтому неявный сигнал вознаграждения хорошо откалиброван. На практике многие наборы данных (например, Anthropic HH-RLHF или сравнения OpenAI WebGPT) сопоставляют ответы, сгенерированные моделью, с ответами человека или ответами более слабой модели\,---\,допустимый подход, хотя он вносит дополнительное несоответствие распределений.

\textbf{Распределение промптов.} Промпты $x$ должны отражать целевое распределение развёртывания. Модель, обученная с DPO, будет хорошо работать на промптах, схожих с теми, что есть в $\mathcal{D}$, и может деградировать на промптах из других распределений\,---\,та же оговорка об обобщении, что применима ко всем методам обучения с учителем.

\subsection{Детали реализации}

\textbf{Вычисление логарифмической вероятности.} Логарифм вероятности последовательности $\log\pi_\theta(y|x) = \sum_{t=1}^{|y|} \log\pi_\theta(y_t | x, y_{<t})$ вычисляется путём конкатенации промпта и ответа, выполнения одного прямого прохода и суммирования по-токенных логарифмических вероятностей на позициях ответа.

\textbf{Нормализация по длине.} Ответы разной длины будут иметь разные логарифмические вероятности просто из-за длины. Чтобы стратегия не выучила тривиальный шорткат по длине\,---\,присваивая более высокую логарифмическую вероятность более коротким ответам вне зависимости от качества,\,---\,стандартной практикой является деление на длину ответа: $\bar\ell = \log\pi_\theta(y|x) / |y|$.

\textbf{Кэширование логитов референса.} Референсная модель $\pi_\mathrm{ref}$ заморожена на протяжении всего обучения. Логарифмические вероятности $\log\pi_\mathrm{ref}(y_w|x)$ и $\log\pi_\mathrm{ref}(y_l|x)$ можно поэтому предварительно вычислить один раз и закэшировать, исключив необходимость запуска прямого прохода референсной модели во время обучения. Это экономит примерно половину вычислений на каждом шаге обучения.

\textbf{Гиперпараметры.} Основные гиперпараметры: $\beta \in [0{,}01, 0{,}5]$ (коэффициент KL, управляющий тем, насколько далеко стратегия отходит от референса; меньшее $\beta$ допускает большее отклонение), скорость обучения (как правило, от $1\times10^{-6}$ до $5\times10^{-6}$ для полного дообучения, на один-два порядка меньше, чем при SFT) и число эпох (обычно 1\,---\,3; больше эпох грозит переобучением к данным предпочтений).

\subsection{Сравнение конвейеров DPO и PPO: диаграмма}

Следующая диаграмма сопоставляет полный конвейер PPO-RLHF (сверху) с конвейером DPO (снизу), иллюстрируя компоненты, которые DPO устраняет.

\begin{figure}[ht]
\centering
\resizebox{\textwidth}{!}{%
\begin{tikzpicture}[
  every node/.style={font=\small},
  box/.style={draw, rounded corners=3pt, minimum width=2.4cm, minimum height=1.3cm,
              align=center, text width=2.3cm, inner sep=2pt},
  eliminated/.style={draw, rounded corners=3pt, minimum width=2.4cm, minimum height=1.3cm,
              align=center, text width=2.3cm, inner sep=2pt,
              fill=red!10, draw=red!50, text=red!70},
  active/.style={draw, rounded corners=3pt, minimum width=2.4cm, minimum height=1.3cm,
              align=center, text width=2.3cm, inner sep=2pt, fill=blue!8},
  arrow/.style={->, >=Stealth, thick},
  rowlabel/.style={font=\small\bfseries}
]
\def\dx{3.2}     % шаг между блоками
\def\xlab{-0.4}  % позиция метки PPO:/DPO: (меньше левое поле)
\def\xstart{2.0} % центр первого блока (больше зазор от метки)
\def\yP{2.6}
\def\yD{-0.6}

%% ---- Строка PPO ----
\node[rowlabel, anchor=east] at (\xlab,\yP) {PPO:};
\node[active] (ppo-ref) at (\xstart+0*\dx,\yP) {Референсная стратегия $\pi_\mathrm{ref}$};
\node[active] (ppo-sft) at (\xstart+1*\dx,\yP) {Роллаут: $y \sim \pi_\theta$};
\node[active] (ppo-rm)  at (\xstart+2*\dx,\yP) {Модель вознаграждения $r_\psi$};
\node[active] (ppo-vf)  at (\xstart+3*\dx,\yP) {Функция ценности $V_\phi$};
\node[active] (ppo-ppo) at (\xstart+4*\dx,\yP) {Обновление PPO};

\draw[arrow] (ppo-ref) -- node[above,font=\scriptsize]{KL}       (ppo-sft);
\draw[arrow] (ppo-sft) -- node[above,font=\scriptsize]{$r(x,y)$} (ppo-rm);
\draw[arrow] (ppo-rm)  -- node[above,font=\scriptsize]{преим.}   (ppo-vf);
\draw[arrow] (ppo-vf)  -- (ppo-ppo);

%% ---- Строка DPO ----
\node[rowlabel, anchor=east] at (\xlab,\yD) {DPO:};
\node[active]     (dpo-ref)  at (\xstart+0*\dx,\yD) {Референсная стратегия $\pi_\mathrm{ref}$};
\node[active]     (dpo-pol)  at (\xstart+1*\dx,\yD) {Стратегия $\pi_\theta$};
\node[eliminated] (dpo-rm)   at (\xstart+2*\dx,\yD) {Модель вознаграждения \xmark};
\node[eliminated] (dpo-vf)   at (\xstart+3*\dx,\yD) {Функция ценности \xmark};
\node[active]     (dpo-loss) at (\xstart+4*\dx,\yD) {Потеря DPO};

\draw[arrow] (dpo-ref) -- node[above,font=\scriptsize]{$\log\pi_\mathrm{ref}$} (dpo-pol);
\draw[arrow] (dpo-pol) -- node[above,font=\scriptsize]{$\log\pi_\theta$}       (dpo-rm);
\draw[arrow,draw=red!40,dashed] (dpo-rm) -- (dpo-vf);
\draw[arrow,draw=red!40,dashed] (dpo-vf) -- (dpo-loss);

% Скобка «исключено в DPO»
\draw[decorate, decoration={brace, amplitude=5pt, mirror}, color=red!60]
  ([yshift=-6pt]dpo-rm.south west) -- ([yshift=-6pt]dpo-vf.south east)
  node[midway, below=7pt, font=\scriptsize, color=red!60]{исключено в DPO};

% Прямая стрелка в обход исключённых компонентов
\draw[arrow, color=blue!70]
  (dpo-pol.south) .. controls +(0,-3.2) and +(0,-3.2) .. (dpo-loss.south)
  node[pos=0.5, below=2pt, font=\scriptsize, color=blue!70]{напрямую (без RM)};

\end{tikzpicture}%
}
\caption{Конвейеры обучения PPO (сверху) и DPO (снизу). Конвейер PPO требует четырёх компонентов: референсной стратегии (для KL), онлайн-роллаутов, модели вознаграждения $r_\psi$ и критика с функцией ценности $V_\phi$. DPO устраняет модель вознаграждения и функцию ценности, вычисляя сигнал выравнивания напрямую из логарифмических отношений вероятностей.}
\label{fig:dpo-vs-ppo}
\end{figure}

% ===========================================================
\section{Варианты DPO}
\label{sec:dpo-variants}

Цель DPO, при всей её элегантности, делает конкретные предположения о модели предпочтений и природе обучающих данных. Для смягчения этих предположений и устранения практических ограничений было предложено несколько вариантов.

\subsection{IPO: оптимизация идентичных предпочтений}
\index{IPO (Identity Preference Optimisation)}

DPO выводится путём максимизации логарифмического правдоподобия в рамках модели Брэдли\,---\,Терри, которая является \emph{мягкой} моделью предпочтений: даже детерминированные предпочтения (один ответ выбирается всегда) дают конечную потерю при любой конечной стратегии. Однако Azar и др. (2023) заметили, что если предпочтения совершенно последовательны\,---\,реалистичная идеализация при использовании автоматического оракула предпочтений,\,---\,оптимум DPO вырождается: стратегия стимулируется увеличивать отступ вознаграждения $h_\theta$ без ограничений, отправляя логарифмические вероятности к $+\infty$ и $-\infty$. Это \emph{переобучение к детерминированным предпочтениям} проявляется как вырожденные выходные распределения.

Функция потерь Identity Preference Optimisation (IPO) устраняет эту проблему, заменяя логарифм сигмоиды штрафом на основе квадрата ошибки:
\begin{equation}
\label{eq:ipo-loss}
\mathcal{L}_\mathrm{IPO}(\theta)
= \frac{1}{N}\sum_{i=1}^{N}
\left(
\log\frac{\pi_\theta(y_w^{(i)}|x^{(i)})}{\pi_\mathrm{ref}(y_w^{(i)}|x^{(i)})}
- \log\frac{\pi_\theta(y_l^{(i)}|x^{(i)})}{\pi_\mathrm{ref}(y_l^{(i)}|x^{(i)})}
- \frac{1}{2\beta}
\right)^2.
\end{equation}
Целевой отступ вознаграждения равен $1/(2\beta)$: стратегия должна выучить присваивать разность неявных вознаграждений ровно $1/(2\beta)$ между предпочтительным и нежелательным ответами, не больше. Квадратичная потеря создаёт восстанавливающую силу, штрафующую как недостаточное, так и избыточное разделение, предотвращая режим неограниченного отступа. IPO можно вывести непосредственно из исходной KL-ограниченной цели, не проходя через модель Брэдли\,---\,Терри, что делает его применимым даже когда модель предпочтений нелогистическая.

\subsection{KTO: оптимизация Канемана\,---\,Тверски}
\index{KTO (Kahneman-Tversky Optimisation)}

Фундаментальным требованием DPO (и IPO) являются \emph{парные} данные предпочтений: для каждого промпта $x$ должен быть доступен как предпочтительный ответ $y_w$, так и нежелательный ответ $y_l$. Сбор таких парных данных требует одновременного предъявления аннотатору двух ответов\,---\,интерфейса сравнения, который дороже и менее естественен, чем простой вопрос о том, хорош ли отдельный ответ.

Ethayarajh и др. (2023) предложили оптимизацию Канемана\,---\,Тверски (KTO), работающую с непарной, бинарной обратной связью\,---\,каждый пример $(x, y, z)$ состоит из промпта, ответа и бинарной метки $z \in \{0, 1\}$, указывающей, хорош ли ответ ($z=1$) или плох ($z=0$). KTO выводит свою цель из \emph{теории перспектив}\index{prospect theory} (Канеман и Тверски, 1979), в частности из наблюдения, что люди не склонны к риску: боль от потери превышает удовольствие от эквивалентного выигрыша.

\begin{definition}[Функция потерь KTO]
\label{def:kto-loss}
\index{KTO loss}
Пусть $\mathcal{D}_+ = \{(x, y) : z=1\}$ и $\mathcal{D}_- = \{(x, y) : z=0\}$. Опорная точка KTO равна
\[
z_\mathrm{ref} = \E_{(x,y) \in \mathcal{D}} \!\left[\beta\,\KL{\pi_\theta(\cdot|x)}{\pi_\mathrm{ref}(\cdot|x)}\right],
\]
на практике аппроксимируемая как $\beta \log\pi_\theta(y|x)/\pi_\mathrm{ref}(y|x)$, усреднённая по батчу. \textbf{Функция потерь KTO} имеет вид
\begin{equation}
\label{eq:kto-loss}
\mathcal{L}_\mathrm{KTO}(\theta)
= \E_{(x,y)\in\mathcal{D}_+}\!\bigl[1 - \sigma(\hat{r}_\theta(x,y) - z_\mathrm{ref})\bigr]
+ \E_{(x,y)\in\mathcal{D}_-}\!\bigl[1 - \sigma(z_\mathrm{ref} - \hat{r}_\theta(x,y))\bigr],
\end{equation}
где $\hat{r}_\theta(x,y) = \beta\log(\pi_\theta(y|x)/\pi_\mathrm{ref}(y|x))$\,---\,неявное вознаграждение.
\end{definition}

Интуитивно, KTO поощряет неявное вознаграждение положительных примеров превышать опорную точку (тогда человек «получает» полезность) и неявное вознаграждение отрицательных примеров падать ниже неё (тогда человек «теряет»\,---\,но фрейм неприятия потерь означает, что это штрафуется асимметрично). KTO особенно привлекателен, когда наборы данных предпочтений собираются через интерфейсы оценки «нравится / не нравится», которые гораздо более распространены в масштабе, чем интерфейсы сравнения двух ответов.

\subsection{cDPO: консервативный DPO}
\index{cDPO (Conservative DPO)}

Аннотации предпочтений людей зашумлены: аннотаторы иногда предпочитают объективно худший ответ из-за поверхностных признаков, нехватки времени или непоследовательности. Стандартный DPO обрабатывает все аннотации с одинаковым весом, что означает: пара с перевёрнутой меткой\,---\,где $y_w$ на самом деле хуже $y_l$,\,---\,будет активно двигать стратегию в неправильном направлении.

Консервативный DPO (cDPO), предложенный Mitchell и др. (2023), устраняет эту проблему, добавляя параметр сглаживания меток $\epsilon \in [0, 0{,}5)$ в цель DPO. Вместо трактовки каждого предпочтения как жёсткого присвоения cDPO трактует его как мягкий сигнал с вероятностью шума $\epsilon$:
\begin{equation}
\label{eq:cdpo-loss}
\mathcal{L}_\mathrm{cDPO}(\theta)
= -\frac{1}{N}\sum_{i=1}^{N}
\Bigl[
(1 - \epsilon)\log\sigma(h^{(i)}_\theta)
+ \epsilon\log\sigma(-h^{(i)}_\theta)
\Bigr],
\end{equation}
где $h^{(i)}_\theta$\,---\,разность неявных вознаграждений, определённая в определении~\ref{def:dpo-loss}. При $\epsilon = 0$ cDPO точно сводится к DPO. При $\epsilon > 0$ функция потерь добавляет небольшой штраф за чрезмерную уверенность в том, что $y_w$ лучше, что регуляризует стратегию против переобучения к зашумлённым меткам. Цель cDPO эквивалентна потере модели вознаграждения со сглаживанием меток~\eqref{eq:rm-label-smooth}, применённой в параметризации DPO.

\subsection{ORPO: оптимизация отношения шансов}
\index{ORPO (Odds Ratio Preference Optimisation)}

Все упомянутые выше варианты отделяют цель выравнивания от цели языкового моделирования: DPO обучается на парах предпочтений, тогда как базовое SFT-обучение является отдельным предварительным этапом. Hong и др. (2024) предложили Odds Ratio Preference Optimisation (ORPO), который объединяет две цели в единую функцию потерь и полностью устраняет необходимость в референсной модели.

\begin{definition}[Функция потерь ORPO]
\label{def:orpo-loss}
\index{ORPO loss}
\textbf{Функция потерь ORPO} объединяет обучающий член отрицательного логарифмического правдоподобия для предпочтительного ответа с штрафом на основе отношения шансов:
\begin{equation}
\label{eq:orpo-loss}
\mathcal{L}_\mathrm{ORPO}(\theta)
= -\log\pi_\theta(y_w|x)
- \lambda\,\log\sigma\!\left(\log\frac{\mathrm{odds}_\theta(y_w|x)}{\mathrm{odds}_\theta(y_l|x)}\right),
\end{equation}
где $\mathrm{odds}_\theta(y|x) = \pi_\theta(y|x) / (1 - \pi_\theta(y|x))$\,---\,шансы ответа при данной стратегии, а $\lambda > 0$\,---\,коэффициент смешивания.
\end{definition}

Первый член обеспечивает продолжение генерации плавного, связного текста (дообучение с учителем на предпочтительном ответе). Второй член штрафует случаи, когда шансы стратегии генерировать $y_l$ сравнимы с шансами генерировать $y_w$ или превышают их. Поскольку ORPO использует \emph{отношение шансов}, а не логарифмическое отношение с участием референсной модели, отдельная референсная модель не нужна: регуляризация встроена в структуру функции потерь.

ORPO особенно привлекателен для очень больших моделей, где поддерживать копию референсной модели в памяти нецелесообразно, а также для случаев, когда модель обучается с нуля на данных предпочтений без предварительного этапа SFT.

% ===========================================================
\section{Ограничения прямого выравнивания}
\label{sec:dpo-limitations}

\subsection{Проблема сдвига распределения}

Фундаментальное ограничение DPO и всех автономных прямых алгоритмов выравнивания\,---\,\emph{сдвиг распределения}\index{distribution shift!DPO}. Набор данных предпочтений $\mathcal{D}$ собирается путём предъявления ответов от некоторой порождающей данные стратегии $\mu$\,---\,как правило, SFT-модели $\pi_\mathrm{ref}$\,---\,аннотаторам-людям. Неявный сигнал вознаграждения в DPO откалиброван под это распределение: логарифмическое отношение $\log(\pi_\theta / \pi_\mathrm{ref})$ мало (а градиент потерь велик) именно для ответов из носителя $\mu$.

По мере обучения $\pi_\theta$ отдаляется от $\pi_\mathrm{ref}$. Ответы, изначально вероятные при $\mu$, становятся всё менее вероятными при $\pi_\theta$, а новые ответы, которые $\pi_\theta$ научилась предпочитать, могут вовсе отсутствовать в $\mathcal{D}$. Функция потерь DPO не может давать никакого обучающего сигнала для ответов вне своего обучающего распределения. Это означает, что по мере улучшения стратегии она исследует области пространства ответов, где имеющиеся у неё метки предпочтений всё менее информативны.

Эмпирически это проявляется в характерной кривой обучения: отступ вознаграждения DPO $h_\theta$ быстро растёт в первые несколько тысяч шагов, пока стратегия усваивает очевидные различия предпочтений в данных, а затем выходит на плато, когда стратегия исчерпывает информационное содержание фиксированного набора данных. Дальнейшее обучение даёт всё меньшую отдачу, а продолжение обучения после этой точки может даже привести к \emph{деградации} стратегии\,---\,увеличению неявного отступа вознаграждения для обучающих пар при ухудшении производительности на отложенных промптах, что является формой переобучения к предпочтениям.

\subsection{Отсутствие Online Exploration}

Связанным с сдвигом распределения является отсутствие online exploration\index{exploration!absence in DPO}. В online-методах RL, таких как PPO, стратегия непрерывно генерирует новые ответы, собирает вознаграждения и обновляется. Это позволяет стратегии исследовать пространство ответов и получать обратную связь именно по тем ответам, которые она с наибольшей вероятностью генерирует в данный момент. Сигнал предпочтений всегда on-policy.

У DPO такого механизма нет. Он работает исключительно с фиксированным, заранее собранным набором данных. Если существуют высококачественные ответы, которые SFT-модель не смогла бы сгенерировать (и, следовательно, они отсутствуют в $\mathcal{D}$), DPO не может их обнаружить. Это особенно существенно для задач, требующих композиционного рассуждения: модель, научившаяся быть чуть более точной в задачах пошагового рассуждения, теперь может выиграть от обратной связи по конкретным паттернам рассуждения, которые она генерирует,\,---\,а они отличаются от тех, что есть в фиксированном обучающем наборе.

\subsection{Поиск моды против покрытия мод}

\index{mode-seeking!DPO}
Модель Брэдли\,---\,Терри и функция потерь DPO выводятся из перспективы \emph{поиска моды}: потеря поощряет стратегию присваивать высокую вероятность предпочтительным ответам и низкую\,---\,нежелательным. В пределе это толкает стратегию к \emph{точечной массе} на ответе с наибольшим вознаграждением для каждого промпта.

Эта тенденция к поиску моды иногда желательна (для задач с единственным правильным ответом, например математических), но нередко нежелательна (для задач, требующих разнообразия, таких как творческое письмо или открытый диалог). Стратегия, ищущая моду, может генерировать ответы с высокими оценками по метрикам предпочтений, при этом будучи монотонной и лишённой разнообразия, которое реально хотят пользователи. В отличие от неё онлайн-методы RL с регуляризацией по энтропии или бонусами за разнообразие можно настроить для получения более широко охватывающих распределений.

\subsection{Когда DPO не работает}

На основе приведённого выше анализа DPO, как правило, уступает онлайн-RL в следующих сценариях.

\begin{itemize}[leftmargin=*]
  \item \textbf{Длительные горизонты обучения}: по мере потребления большего количества данных предпочтений и расхождения стратегии с референсом накапливается сдвиг распределения, и производительность выходит на плато или регрессирует.
  \item \textbf{Задачи с проверяемой правильностью}: для математических рассуждений или кодирования онлайн-верификатор (юнит-тесты, символический проверщик) обеспечивает более сильный и надёжный сигнал, чем статичные аннотации предпочтений людей.
  \item \textbf{Выравнивание с высокими ставками}: когда цель состоит в создании модели, безопасной и полезной \emph{по всему распределению развёртывания}, онлайн-обратная связь от реальных пользователей даёт сигнал, который фиксированный набор данных не может предвосхитить.
  \item \textbf{Дальнейшее улучшение}: DPO сходится к оптимальной стратегии для данного фиксированного набора данных. Дальнейшее улучшение требует новых данных, тогда как онлайн-RL может продолжать улучшаться, пока среда обеспечивает обратную связь.
\end{itemize}

% ===========================================================
\section{Online RL в Production-системах}
\label{sec:online-rl-production}

\subsection{Почему Online RL важен}

\index{online RL!production systems}
Различие между автономным и онлайн-выравниванием не просто академическое. Производственная языковая модель обслуживает миллионы пользователей, генерируя миллиарды ответов по огромному разнообразию тем, стилей и намерений. Ни один статичный набор данных предпочтений, сколь бы он ни был велик, не может предвидеть это разнообразие. Пользователи будут задавать вопросы, о которых создатели набора данных не задумывались, запрашивать стили, отличающиеся от обучающего распределения, и обнаруживать слабые места, проявляющиеся только в масштабе.

Онлайн-обучение с подкреплением замыкает этот круг. Вместо однократного обучения на фиксированном наборе данных и развёртывания онлайн-система RL непрерывно собирает обратную связь от реальных взаимодействий\,---\,явную (оценки «нравится / не нравится», звёздочки, выборы предпочтений) или неявную (продолжительность сессии, уточняющие вопросы, частота копирования)\,---\,и обновляет стратегию. Поведение модели улучшается по мере обслуживания пользователей, причём улучшение направлено на реальное распределение развёртывания, а не на прокси-набор данных.

Именно такой паттерн развёртывания присущ ведущим коммерческим ИИ-ассистентам. Технический отчёт GPT-4 (OpenAI, 2023) описывает непрерывные обновления RLHF на основе обратной связи пользователей. Конвейер Constitutional AI от Anthropic (Bai и др., 2022) использует RLHF с сэмплированием на текущей стратегии, чтобы сигнал вознаграждения был откалиброван под распределение выходных данных текущей стратегии. Эти системы не рассматривают выравнивание как однократное событие обучения; они рассматривают его как непрерывный производственный процесс.

\subsection{On-Policy и Off-Policy Соображения}

\index{on-policy!alignment}
В online-режиме ключевым проектным решением является использование \emph{on-policy feedback} или \emph{off-policy feedback}. On-policy feedback собирается от текущей стратегии: модель генерирует ответ, пользователь даёт обратную связь, и модель обновляется на этой обратной связи перед генерацией следующего ответа. Off-policy feedback собирается от более ранней версии стратегии (или другой стратегии) и используется для обновления текущей стратегии.

On-policy feedback\,---\,золотой стандарт для выравнивания. Гарантируется, что сигнал вознаграждения находится в распределении, которое стратегия в данный момент исследует, поэтому нет сдвига распределения. Обновление эквивалентно policy gradient с выборочной оценкой advantage. Однако on-policy feedback медленна: обратную связь необходимо собрать и обработать до следующего обновления, что ограничивает пропускную способность цикла обучения.

Off-policy feedback (например, DPO на сохранённых данных предпочтений) обеспечивает значительно более высокую пропускную способность, поскольку обучение отделено от сбора данных. Платой является сдвиг распределения, рассмотренный в разделе~\ref{sec:dpo-limitations}. На практике production-системы обычно используют гибридный подход: большой replay buffer off-policy обратной связи из исторических данных в сочетании с небольшим потоком on-policy feedback, с importance weighting для коррекции несоответствия распределений.

\subsection{Смягчение катастрофического забывания}

\index{catastrophic forgetting}
Постоянной проблемой онлайн-RL является \emph{катастрофическое забывание}: по мере обновления стратегии для улучшения выравнивания она может деградировать по способностям, присутствовавшим в SFT-модели. Модель может стать менее беглой, менее фактически точной или менее способной к задачам, не представленным в текущем потоке обратной связи.

Несколько техник смягчают катастрофическое забывание в производстве.

\begin{description}[leftmargin=!,labelwidth=3.5cm]
  \item[Якорение KL] Штраф KL $\beta\KL{\pi_\theta}{\pi_\mathrm{ref}}$ в цели RL не позволяет стратегии слишком далеко отходить от контрольной точки SFT. Однако если сама контрольная точка SFT периодически обновляется (что обычно для производства), якорь смещается, потенциально допуская медленный, но неограниченный дрейф.
  \item[Ограничения на способности] После каждого обновления RL запускаются отдельные оценки на отложенных бенчмарках способностей (MMLU, HumanEval, GSM8K). Если какая-либо способность падает ниже порога, обновление откатывается или коэффициент KL увеличивается. Это требует специальной инфраструктуры оценки, но является наиболее принципиальной защитой.
  \item[Буферы воспроизведения] Смешивание старых примеров SFT с градиентом онлайн-RL гарантирует, что стратегия продолжает получать контролируемый сигнал по полному распределению способностей. Это аналогично воспроизведению опыта в DQN: примеры SFT играют роль «якорей» в пространстве опыта.
  \item[Дообучение на основе LoRA] Ограничение обновлений адаптером низкого ранга (LoRA) вместо полных параметров модели ограничивает степень забывания, поскольку веса базовой модели остаются фиксированными. Выравнивание хранится в адаптере; способности сохраняются в базовой модели.
\end{description}

\subsection{Ограничения безопасности в онлайн-RL}

\index{safety constraints!online RL}
Онлайн-RL вносит риск безопасности, которого нет у автономного выравнивания: модель обновляется на обратной связи реальных пользователей, которая может включать состязательные, предвзятые или иным образом вредоносные сигналы. Производственная система должна защищаться от нескольких режимов отказа.

\textbf{Взлом вознаграждения пользователями.} Если пользователи обнаружат, что определённые паттерны ответов получают высокие вознаграждения (например, ответы, соглашающиеся с любой позицией пользователя, или особенно подробные ответы), они могут использовать это для управления поведением модели. Состязательная устойчивость сигнала вознаграждения требует обнаружения аномалий и фильтрации выбросов обратной связи.

\textbf{Нестабильность обратной связи.} Положительные обратные связи могут вызвать быстрый дрейф стратегии к локальной моде. Штраф KL в сочетании с периодической оценкой на разнообразном отложенном наборе промптов является стандартным стабилизатором.

\textbf{Оптимизация стратегии с ограничениями.} Для критически важных с точки зрения безопасности свойств (отказ от вредоносных запросов, фактическая точность) можно дополнить цель RL ограничениями Лагранжа (Achiam и др., 2017; Bai и др., 2022): стратегия оптимизирует вознаграждение при условии, что релевантные безопасности метрики остаются выше нижней границы. Это преобразует неограниченную задачу RL в задачу ограниченной оптимизации, решаемую методами двойственных переменных.

\subsection{PPO против DPO в онлайн-режиме}

\index{PPO!vs DPO online}
Сравнение PPO и DPO качественно различается в онлайн- и автономном режимах. В автономном режиме (фиксированный набор данных предпочтений) DPO конкурентоспособен с PPO или превосходит его по многим бенчмаркам: он проще, стабильнее и требует меньше вычислений. В онлайн-режиме PPO имеет структурное преимущество, которое DPO не может компенсировать.

PPO собирает обратную связь по текущей стратегии, обновляет стратегию на основе этой обратной связи и повторяет. Стратегия и источник обратной связи всегда синхронизированы. DPO, напротив, требует отдельного этапа сбора данных перед каждым этапом обучения; каждый цикл «сбор\,---\,обучение» является отдельным батчем, а не непрерывным циклом. Это означает, что онлайн-DPO фактически является \emph{итеративным DPO}\index{iterative DPO}: сбор предпочтений от текущей стратегии, запуск DPO на несколько эпох, сбор новых предпочтений от обновлённой стратегии и повторение. Итеративный DPO существенно сокращает разрыв с PPO за счёт уменьшения сдвига распределения, но вносит свои инженерные издержки (повторный сбор данных и аннотирование) и всё равно уступает полностью онлайн-циклу PPO.

На практике выбор между PPO и DPO в production зависит от доступной инфраструктуры. Команда с выделенной RL-инфраструктурой, быстрой моделью вознаграждения (или верификатором) и возможностью online rollouts должна предпочесть PPO или вариант PPO (например, GRPO). Команда без RL-инфраструктуры или нуждающаяся в быстром выравнивании модели с имеющимся набором данных предпочтений должна предпочесть DPO или один из его вариантов.

\subsection{Практические паттерны развёртывания}

На основе приведённого анализа и наблюдаемой практики в ведущих ИИ-лабораториях сложились следующие паттерны развёртывания.

\begin{description}[leftmargin=!,labelwidth=3.5cm]
  \item[Начальный запуск с DPO] Использовать DPO на имеющемся наборе данных предпочтений для быстрого создания выровненной модели. Это самый быстрый путь от SFT к развёртываемой выровненной модели, подходящий для первоначальных выпусков.
  \item[Уточнение с онлайн-RL] После развёртывания собирать обратную связь по текущей стратегии и запускать итеративный DPO или PPO для непрерывного улучшения. Модель, запущенная с DPO, обеспечивает хорошую отправную точку с относительно небольшим доступным бюджетом KL для дальнейшего улучшения.
  \item[Разделение безопасности и способностей] Использовать автономный DPO для выравнивания безопасности (которое меняется медленно и может быть закодировано в статичном наборе данных) и онлайн-RL для выравнивания полезности (которое должно адаптироваться к меняющемуся распределению развёртывания).
  \item[RL с верификатором для рассуждений] Для задач с проверяемой правильностью (математика, код) использовать онлайн-RL с автоматическим верификатором в качестве сигнала вознаграждения, а не предпочтения людей. Это позволяет неограниченный сбор данных по текущей стратегии без затрат на аннотирование людьми, а верификатор свободен от предвзятостей, влияющих на аннотаторов-людей.
\end{description}

% ===========================================================
\section{Исчерпывающее сравнение методов выравнивания}
\label{sec:alignment-comparison}

\subsection{Обзор методов}

Таблица~\ref{tab:alignment-comparison} содержит исчерпывающее сравнение основных методов выравнивания для больших языковых моделей. Методы охватывают весь спектр\,---\,от полностью онлайн-методов градиента стратегии до чисто автономного выравнивания с учителем.

\begin{table}[ht]
\centering
\caption{Исчерпывающее сравнение методов выравнивания LLM. «RM» = модель вознаграждения; «реф.\ модель» = требуется референсная стратегия; «парные» = требуются парные данные $(y_w, y_l)$; «онлайн» = генерирует новые ответы во время обучения. Объём памяти относительно одной копии модели ($1\times$). Стабильность: В = высокая, С = средняя, Н = низкая.}
\label{tab:alignment-comparison}
\footnotesize
\begin{tabular}{@{}lcccccp{2.2cm}@{}}
\toprule
\textbf{Метод} & \textbf{RM?} & \textbf{Парные?} & \textbf{Онлайн?} & \textbf{Память} & \textbf{Стаб.} & \textbf{Лучше всего для} \\
\midrule
PPO        & Да  & Нет & Да  & $4\times$ & С & Безопасность, RLHF в масштабе, онлайн-обратная связь \\
REINFORCE  & Да  & Нет & Да  & $2\times$ & Н & Простые задачи, базовое сравнение \\
RLOO       & Да  & Нет & Да  & $2\times$ & С & Онлайн-RL с малыми ресурсами, без критика \\
GRPO       & Нет* & Нет & Да & $2\times$ & В & Рассуждения, проверяемые задачи, без критика \\
DPO        & Нет & Да  & Нет & $2\times$ & В & Быстрое выравнивание, автономные данные \\
IPO        & Нет & Да  & Нет & $2\times$ & В & Детерминированные / автоматические предпочтения \\
KTO        & Нет & Нет & Нет & $2\times$ & В & Бинарная обратная связь, без попарного сравнения \\
cDPO       & Нет & Да  & Нет & $2\times$ & В & Зашумлённые метки предпочтений \\
ORPO       & Нет & Да  & Нет & $1\times$ & В & Без реф.\ модели, ограниченная память \\
\bottomrule
\end{tabular}

\smallskip
{\footnotesize *GRPO использует верификатор или вознаграждение на основе правил, а не обученную модель вознаграждения.}
\end{table}

\begin{table}[ht]
\centering
\caption{Плюсы и минусы каждого метода выравнивания.}
\label{tab:alignment-pros-cons}
\footnotesize
\begin{tabular}{@{}lp{3.5cm}p{3.5cm}@{}}
\toprule
\textbf{Метод} & \textbf{Плюсы} & \textbf{Минусы} \\
\midrule
PPO
  & Полностью онлайн, обрабатывает сложные сигналы вознаграждения, хорошо изучен теоретически
  & Высокое потребление памяти (4 модели), сложная реализация, риск взлома вознаграждения \\
\midrule
REINFORCE
  & Простота, несмещённый градиент стратегии, лёгкость реализации
  & Высокая дисперсия, медленная сходимость, нестабильность без baseline \\
\midrule
RLOO
  & Снижение дисперсии через базис leave-one-out, не требует критика
  & Требует группового сэмплирования, по-прежнему имеет сложность онлайн-режима \\
\midrule
GRPO
  & Нет критика, стабильный, отлично работает на проверяемых задачах, сильные эмпирические результаты
  & Требует верификатора или модели вознаграждения; нет естественного расширения на открытые задачи \\
\midrule
DPO
  & Простота, стабильность, обучение в стиле SFT, нет RM, низкое потребление памяти
  & Только автономный, сдвиг распределения, нет исследования, выходит на плато \\
\midrule
IPO
  & Устраняет переобучение DPO, регулирует отступ, более устойчив к чистым предпочтениям
  & Чувствителен к гиперпараметрам ($\beta$ управляет целевым отступом), нет сигнала по текущей стратегии \\
\midrule
KTO
  & Работает с бинарными метками (без пар), естественен для обратной связи в производстве
  & Требует тщательной оценки опорной точки, менее эффективен по данным, чем DPO \\
\midrule
cDPO
  & Устойчив к шуму меток, естественный регуляризатор, почти нулевые накладные расходы по сравнению с DPO
  & Требует настройки $\epsilon$; слишком большое $\epsilon$ препятствует выравниванию \\
\midrule
ORPO
  & Не нужна референсная модель, SFT + выравнивание за один проход, экономия памяти
  & Меньше теоретического обоснования; $\lambda$ требует настройки; нет регуляризации KL \\
\bottomrule
\end{tabular}
\end{table}

\subsection{Выбор метода выравнивания по цели}

Выбор метода выравнивания должен определяться целью развёртывания, имеющимися данными и доступной вычислительной инфраструктурой. Следующие рекомендации сводят приведённый выше анализ к практическим указаниям.

\textbf{Выравнивание безопасности.} Для встраивания ограничений безопасности\,---\,отказа от вредоносных запросов, избегания токсичных выходных данных, соблюдения конституциональных принципов\,---\,основными проблемами являются покрытие (модель должна вести себя безопасно во всём распределении развёртывания) и стабильность (ограничения не должны деградировать со временем). PPO с моделью вознаграждения, специфичной для безопасности, и жёсткими ограничениями Лагранжа является наиболее принципиальным подходом для флагманских моделей. DPO или cDPO на тщательно подобранном наборе данных предпочтений безопасности является практической и нередко достаточной альтернативой для меньших моделей или команд без инфраструктуры RL. KTO особенно хорошо подходит для выравнивания безопасности, поскольку метки безопасности по природе бинарны (вредоносный / не вредоносный), а не сравнительны.

\textbf{Способность к рассуждению.} Для улучшения производительности на математических, логических задачах или задачах кодирования ключевое наблюдение состоит в том, что у этих задач есть проверяемые ответы. Это означает, что вместо аннотаций предпочтений людей можно использовать сигнал вознаграждения на основе правил (правильно / неправильно), обеспечивая неограниченный сбор данных по текущей стратегии. GRPO\index{GRPO (Group Relative Policy Optimisation)} является наиболее сильным выбором здесь: он использует групповое сэмплирование для оценки преимуществ без критика, а вознаграждение на основе верификатора даёт чистый, однозначный сигнал. Результат DeepSeek-R1 (DeepSeek-AI, 2025) продемонстрировал, что GRPO с вознаграждением на основе верификатора может давать драматические улучшения способности к рассуждению без каких-либо данных предпочтений людей.

\textbf{Следование инструкциям.} Следование инструкциям\,---\,способность точно и полезно выполнять разнообразные запросы пользователей\,---\,требует как широкого охвата (множество типов инструкций), так и тонких оценок качества (разница между хорошим и отличным ответом на данную инструкцию). DPO на большом разнообразном наборе данных предпочтений (например, UltraFeedback или Alpaca Farm) является рабочим методом для следования инструкциям. Итеративный DPO\,---\,несколько раундов DPO с сбором данных по текущей стратегии между раундами\,---\,дополнительно улучшает охват. PPO с обобщённой моделью вознаграждения является наиболее сильным вариантом при наличии онлайн-обратной связи.

\textbf{Творческие задачи.} Творческие задачи (генерация историй, поэзия, открытый диалог) требуют разнообразия: выровненная модель должна производить разнообразные, образные выходные данные, а не сходиться к единой моде. Тенденция DPO к поиску моды делает его неподходящим выбором, если только не добавить регуляризацию по разнообразию. Предпочтительнее PPO с бонусом за энтропию или функцией вознаграждения, поощряющей разнообразие. KTO с положительными метками для разнообразных творческих примеров также может стимулировать широкий охват без тенденции к коллапсу к моде, присущей методам на основе контрастных пар.

\textbf{Среды с ограниченными ресурсами.} При жёстких ограничениях по памяти или вычислениям единственными осуществимыми вариантами являются ORPO (требует только одну копию модели) или DPO (две копии). Оба могут применяться с адаптерами LoRA для дальнейшего сокращения числа обучаемых параметров, позволяя выравнивать большие модели на обычном оборудовании.

\begin{example}[Выбор метода для математического репетиторского ассистента]
\label{ex:alignment-choice}
Команда разрабатывает математического репетиторского ассистента, который должен: (a) правильно решать задачи, (b) отказываться давать прямые ответы, когда студент должен учиться самостоятельно, и (c) ясно объяснять рассуждения. Как следует подходить к выравниванию?

Для цели~(a) идеально подходит GRPO с символическим верификатором математики: ответы можно верифицировать автоматически, данные по текущей стратегии не ограничены, и аннотирование людьми не требуется.

Для цели~(b) небольшой набор данных предпочтений безопасности из пар (задача, «дать ответ», «направлять без ответа») можно выровнять с DPO или KTO. Ограничение похоже на ограничение безопасности (бинарное: нарушает политику «без прямого ответа» или нет), что делает KTO особенно естественным.

Для цели~(c) подходит DPO на наборе данных человеческих пар предпочтений (хорошо объяснённые vs плохо объяснённые решения), дополненный раундами итеративного DPO с использованием роллаутов по текущей стратегии от модели, улучшенной GRPO, для избегания сдвига распределения.

На практике эти три цели объединяются: GRPO обрабатывает цель проверяемой правильности, KTO обрабатывает педагогическое ограничение в стиле безопасности, а DPO обрабатывает стилистическое измерение качества. Каждый метод применяется на отдельном этапе, с оценкой способностей между этапами для мониторинга забывания.
\end{example}

% ===========================================================
\section{Итоги}
\label{sec:ch12-summary}

В этой главе была развита теория и практика прямой оптимизации предпочтений и её связь с широким ландшафтом методов выравнивания LLM.

Центральное понимание DPO (раздел~\ref{sec:dpo-derivation}) состоит в двухшаговой алгебраической манипуляции: оптимальная стратегия KL-регуляризованного RLHF может быть преобразована для выражения вознаграждения как логарифмического отношения оптимальной стратегии к референсной плюс зависящая от промпта статистическая сумма. Когда это выражение подставляется в вероятность предпочтения Брэдли\,---\,Терри, статистическая сумма сокращается, оставляя функцию потерь, зависящую только от логарифмических отношений стратегий\,---\,величин, вычислимых двумя прямыми проходами. Модель вознаграждения устраняется без приближений.

Алгоритм DPO (раздел~\ref{sec:dpo-algorithm}) сводится к бинарной кросс-энтропии на логарифмических отношениях стратегий с весом доверия $\sigma(-h_\theta)$, фокусирующим обучение на неопределённых парах. Диаграмма TikZ (рисунок~\ref{fig:dpo-vs-ppo}) проиллюстрировала драматическое упрощение по сравнению с PPO: модель вознаграждения и функция ценности устранены, уменьшая требуемые копии моделей с четырёх до двух.

Варианты DPO (раздел~\ref{sec:dpo-variants}) каждый устраняют конкретное ограничение. IPO предотвращает режим неограниченного отступа, налагая квадратичную цель на разрыв вознаграждения. KTO расширяет выравнивание на непарную бинарную обратную связь с использованием теории перспектив. cDPO вводит сглаживание меток для обработки зашумлённых аннотаций. ORPO полностью устраняет референсную модель, используя штраф на основе отношения шансов, встроенный в контролируемую функцию потерь.

Проблема сдвига распределения (раздел~\ref{sec:dpo-limitations}) является фундаментальным ограничением, общим для всех offline direct alignment methods. По мере удаления стратегии от reference policy метки предпочтений в обучающем наборе становятся всё менее информативными, вызывая выход производительности на плато. Отсутствие online exploration означает, что DPO не может открыть высококачественные ответы, которые текущая стратегия теперь способна генерировать, но которых не было в исходном наборе данных.

Online RL в production (раздел~\ref{sec:online-rl-production}) устраняет эти ограничения, замыкая контур между поведением модели и обратной связью. Production-система RL непрерывно генерирует ответы, собирает обратную связь (явную или неявную) и обновляет стратегию. Основные инженерные задачи\,---\,катастрофическое забывание (смягчаемое якорением KL, ограничениями на способности, replay buffer и LoRA), безопасность при состязательной обратной связи (смягчаемая обнаружением аномалий и constrained optimisation) и компромисс PPO vs.\ DPO (разрешаемый доступностью инфраструктуры и допустимой степенью сдвига распределения).

Исчерпывающая таблица сравнения (раздел~\ref{sec:alignment-comparison}) является практическим справочником для выбора метода выравнивания по цели: PPO и GRPO для онлайн-режима и проверяемых задач, DPO и его варианты для автономного режима с доступными данными предпочтений, KTO для бинарной обратной связи и ORPO для сред с ограниченной памятью.

\begin{description}[leftmargin=!,labelwidth=3.5cm]
  \item[Глава~\ref{ch:reward-models}] закладывает основу: модель Брэдли\,---\,Терри, KL-регуляризованную цель и параметризацию оптимальной стратегии, которую использует DPO.
  \item[Глава~\ref{ch:practical-rlhf}] описывает конвейер дообучения PPO во всех инженерных деталях, включая вознаграждение KL на уровне токенов, оценку преимущества и практические трюки для стабильного обучения.
  \item[Глава~\ref{ch:agents}] обсуждает выравнивание в агентской среде, где модель должна быть выровнена не только по однооборотным ответам, но и по многошаговым траекториям.
\end{description}

% ===========================================================
\section*{Упражнения}
\addcontentsline{toc}{section}{Упражнения}

\begin{exercise}
\label{ex:dpo-derivation}
Проведите полный вывод DPO с нуля. Начиная с KL-регуляризованной цели~\eqref{eq:dpo-rlhf-obj}, докажите, что оптимальная стратегия имеет вид~\eqref{eq:dpo-optimal-policy}. Затем преобразуйте для получения параметризации вознаграждения~\eqref{eq:dpo-reparam}. Покажите явно, что статистическая сумма $Z(x)$ сокращается при вычислении разности вознаграждений $r(x, y_w) - r(x, y_l)$. Наконец, подставьте в формулу Брэдли\,---\,Терри для получения функции потерь DPO~\eqref{eq:dpo-loss}. На каждом шаге укажите, какие предположения являются точными и что потребовалось бы изменить, если бы модель предпочтений не была Брэдли\,---\,Терри.
\end{exercise}

\begin{exercise}
\label{ex:dpo-gradient}
Проверьте утверждение~\ref{prop:dpo-gradient} явным вычислением. Для одной пары предпочтений с логарифмическими вероятностями $\log\pi_\theta(y_w|x) = -2{,}0$, $\log\pi_\mathrm{ref}(y_w|x) = -2{,}4$, $\log\pi_\theta(y_l|x) = -1{,}5$, $\log\pi_\mathrm{ref}(y_l|x) = -1{,}3$ и $\beta = 0{,}1$:
\begin{enumerate}
  \item Вычислите разность неявных вознаграждений $h_\theta$.
  \item Вычислите потерю DPO для этой пары.
  \item Вычислите вес доверия $\sigma(-h_\theta)$ и интерпретируйте его: является ли данная пара лёгкой или сложной для текущей стратегии?
  \item Предположим, что $\log\pi_\theta(y_w|x)$ уменьшается на $0{,}3$ (стратегия становится менее вероятной для генерации $y_w$). Пересчитайте $h_\theta$ и потерю и объясните качественно, как изменяется градиент.
\end{enumerate}
\end{exercise}

\begin{exercise}
\label{ex:ipo-vs-dpo}
Функция потерь IPO~\eqref{eq:ipo-loss} заменяет логарифм сигмоиды квадратичной потерей.
\begin{enumerate}
  \item Покажите, что функция потерь DPO не имеет конечного минимума, когда предпочтения совершенно последовательны (т.е. когда $y_w \succ y_l$ с вероятностью~$1$). \emph{Подсказка:} рассмотрите, что происходит с $\mathcal{L}_\mathrm{DPO}$ при $h_\theta \to +\infty$.
  \item Покажите, что функция потерь IPO~\eqref{eq:ipo-loss} минимизируется при $h_\theta^* = 1/(2\beta)$. Вычислите $h_\theta^*$ для $\beta = 0{,}1$ и интерпретируйте это значение в терминах разрыва неявного вознаграждения.
  \item Для численного примера из упражнения~\ref{ex:dpo-gradient} вычислите потерю IPO и сравните её с потерей DPO. Какая больше и почему?
\end{enumerate}
\end{exercise}

\begin{exercise}
\label{ex:kto-prospect}
Цель KTO мотивирована теорией перспектив. Формально функция ценности Канемана\,---\,Тверски $v : \R \to \R$ вогнута для выигрышей ($v > 0$) и выпукла и круче для потерь ($v < 0$), с $v(0) = 0$.
\begin{enumerate}
  \item Нарисуйте функцию ценности и объясните, почему неприятие потерь подразумевает $|v'(-\epsilon)| > |v'(\epsilon)|$ для малых $\epsilon > 0$.
  \item Опорная точка KTO $z_\mathrm{ref}$ играет роль нуля функции ценности. Объясните интуитивно, почему неявное вознаграждение положительного примера должно превышать $z_\mathrm{ref}$, а неявное вознаграждение отрицательного примера должно быть ниже него.
  \item Покажите, что KTO с $z_\mathrm{ref} = 0$ и равными долями положительных и отрицательных примеров сводится к форме, близко связанной с DPO, если положительные и отрицательные примеры взяты из одного парного набора данных.
\end{enumerate}
\end{exercise}

\begin{exercise}
\label{ex:cdpo-label-smooth}
Покажите, что функция потерь cDPO~\eqref{eq:cdpo-loss} может быть выведена из DPO, предположив, что каждая метка предпочтения переворачивается с вероятностью $\epsilon$. То есть аннотатор маркирует $y_w \succ y_l$ с вероятностью $1-\epsilon$ и $y_l \succ y_w$ с вероятностью $\epsilon$. Взятие математического ожидания потери DPO по данной стохастической модели маркировки даёт cDPO. Какому значению $\epsilon$ соответствует полностью неинформативный аннотатор и к чему cDPO сводится в этом пределе?
\end{exercise}

\begin{exercise}
\label{ex:distribution-shift}
Рассмотрите стилизованную модель сдвига распределения в DPO. Предположим, набор данных предпочтений собран из референсной стратегии $\pi_\mathrm{ref}$, и после обучения DPO стратегия $\pi_\theta$ имеет дивергенцию KL $d = \KL{\pi_\theta}{\pi_\mathrm{ref}}$ от референса.
\begin{enumerate}
  \item Докажите, что доля ответов в $\pi_\theta$, которые также появляются в носителе $\pi_\mathrm{ref}$ (и, следовательно, покрыты набором данных предпочтений), уменьшается с ростом $d$.
  \item В худшем случае, если $\pi_\theta$ размещает всю свою массу на ответах, не входящих в носитель $\pi_\mathrm{ref}$, что может сообщить стратегии о качестве этих ответов градиент DPO?
  \item Предложите практическое смягчение: как итеративный DPO снижает тяжесть сдвига распределения? Каков остаточный сдвиг распределения после $k$ раундов итеративного DPO?
\end{enumerate}
\end{exercise}

\begin{exercise}
\label{ex:online-rl-design}
Команда развёртывает чат-бот для обслуживания клиентов. Чат-бот должен: (a) правильно решать проблемы клиентов, (b) отвечать вежливо даже на оскорбительные сообщения и (c) переводить разговор к человеку-оператору при необходимости. Разработайте конвейер обучения с выравниванием для этого чат-бота, ответив на следующие вопросы:
\begin{enumerate}
  \item Какой метод(ы) выравнивания вы бы использовали для каждой из трёх целей (a), (b), (c)?
  \item Как бы вы собирали обратную связь в производстве: явную, неявную или обе?
  \item Как бы вы предотвратили деградацию беглости языка или фактической точности модели в цикле онлайн-RL?
  \item Какой механизм безопасности предотвратил бы систематическое управление стратегией со стороны состязательных пользователей?
\end{enumerate}
\end{exercise}

\begin{exercise}
\label{ex:comparison-table}
Изучите сравнение методов выравнивания в таблице~\ref{tab:alignment-comparison}.
\begin{enumerate}
  \item Объясните, почему PPO требует $4\times$ памяти одной модели, тогда как DPO требует только $2\times$.
  \item У команды есть набор данных из 50\,000 взаимодействий с поддержкой клиентов, каждое маркировано как «решено» или «не решено», без попарных сравнений ответов. Какой метод выравнивания вы бы рекомендовали и почему?
  \item GRPO указан с оценкой «В» (высокая) стабильности, несмотря на то что это онлайн-метод. Объясните, почему GRPO достигает более высокой стабильности, чем стандартный REINFORCE, обратившись к механизму групповой нормализации.
  \item При каких условиях DPO мог бы \emph{превзойти} PPO даже в онлайн-режиме? Будьте конкретны в отношении предположений о данных и инфраструктуре.
\end{enumerate}
\end{exercise}

\begin{exercise}
\label{ex:dpo-math-reasoning}
Рассмотрите использование DPO для улучшения математических рассуждений языковой модели. Набор данных предпочтений состоит из пар $(y_w, y_l)$, где $y_w$\,---\,правильное решение с цепочкой мыслей, а $y_l$\,---\,неправильное, оба сгенерированы SFT-моделью.
\begin{enumerate}
  \item После обучения DPO точность модели на математических задачах из обучающего распределения улучшается с 45\% до 58\%. Предложите эксперимент, позволяющий определить, связано ли это улучшение с (i) генерацией лучших шагов рассуждения или (ii) улучшением угадывания формата итогового ответа.
  \item Объясните, почему проблема сдвига распределения (раздел~\ref{sec:dpo-limitations}) особенно тяжела для математических рассуждений по сравнению, скажем, со следованием инструкциям.
  \item Опишите, как вы бы перешли от выравнивания математики на основе DPO к выравниванию на основе GRPO. Как бы выглядела последовательность перехода и что бы вы использовали в качестве сигнала вознаграждения GRPO?
\end{enumerate}
\end{exercise}
